% answers_short.tex —— 2007 年·普通高等学校招生考试·四川卷·理科 · 简短答案速查
% 由 tools/make_short_answers.py 自 answers.tex 的 \gksolhead/\gkanswer 提取，
% 并与 2007/2007_数学真题_四川理_答案要点.md 逐题比对后生成，勿手改（改源 MD 后重跑）。

\begin{gknotice}
\textbf{使用说明}：本节只给最终答案便于快速核对，完整推导见解析版。答案为编写者独立推导，\textbf{非官方标准答案}。
\end{gknotice}

\begin{gknotice}
选择题答案中的 U 表示待核，未给出猜测答案，详情见解析版。
\end{gknotice}

\gkpart{选择题}
\begin{center}\begin{tabular}{ccccccccccccc}
\toprule
题号 & 1 & 2 & 3 & 4 & 5 & 6 & 7 & 8 & 9 & 10 & 11 & 12 \\
\midrule
答案 & A & C & D & D & A & C & A & C & B & B & U & B \\
\bottomrule
\end{tabular}\end{center}

\gkpart{填空题}
\begin{itemize}[leftmargin=2.6em]
  \item[\textbf{13.}] $1$。
  \item[\textbf{14.}] $\frac{\pi}{6}$。
  \item[\textbf{15.}] $x=\frac32$。
  \item[\textbf{16.}] ①④。
\end{itemize}

\gkpart{解答题}
\begin{itemize}[leftmargin=2.6em]
  \item[\textbf{17.}] （1）$\tan2\alpha=-\frac{8\sqrt3}{47}$；（2）$\beta=\frac{\pi}{3}$。
  \item[\textbf{18.}] （1）$0.9984$；（2）$P(\xi=0)=\frac{68}{95}$，$P(\xi=1)=\frac{51}{190}$，$P(\xi=2)=\frac{3}{190}$，$E\xi=\frac{3}{10}$，拒收概率 $\frac{27}{95}$。
  \item[\textbf{19.}] （1）证明见解析（$PC\perp$ 平面 $ABC$ 从而平面 $PAC\perp$ 平面 $ABC$）；（2）$\arctan\frac{2\sqrt3}{3}$，约 $49.1^\circ$；（3）$\frac{\sqrt3}{12}$。
  \item[\textbf{20.}] （1）最大值 $1$，最小值 $-2$；（2）$k\in\left(-2,-\frac{\sqrt3}{2}\right)\cup\left(\frac{\sqrt3}{2},2\right)$。
  \item[\textbf{21.}] （1）$x_{n+1}=\frac{x_n^2+4}{2x_n}$；（2）充要条件是 $x_1\geqslant2$（证明见解析）；（3）$\{a_n\}$ 是首项 $\lg3$、公比 $2$ 的等比数列，$x_n=\frac{2\left(3^{2^{n-1}}+1\right)}{3^{2^{n-1}}-1}$。
  \item[\textbf{22.}] （1）$\frac{20}{n^3}$；（2）证明见解析；（3）在题设 $n>1$ 即 $n\geqslant2$ 上存在且唯一，$a=2$；若把 $n=1$ 计入，则 $S_1=2$ 使严格不等式无整数解。
\end{itemize}
